\section{Notations}
\label{s:nota}

In the paper, we use various LTSs and behavioural relations for them, 
both for OGS and for the $\pi$-calculus.  
In this section, we introduce or summarise common notations. 

We use a tilde, like in $\til a$, for  (possibly empty) tuples of objects (usually names). 
Let $ K \arrG \mu K'$ be a generic LTS (for OGS or \piI; the grammar for actions in the
LTSs  for OGS
and \piI will be the same). 
Actions, ranged over by $\mu$,  can be of the form $\inp a \tilb$, $\bout a \tilb$,
$\tau$,  and $(\tila)$, where $\tau$, called \emph{silent} or (\emph{invisible}) action, 
represents an internal step in $K$, that is, an action that does not require interaction
with the outside,  and $(\tila)$ is  a special action performed by abstractions in \piI
and initial configurations in OGS.   If $\mu \neq \tau$ then $\mu$ is a \emph{visible} action; we use
$\ell$ to range over them.

We sometimes abbreviate $\arrG\tau$ as $ \longrightarrowG $.
We  write $ \LongrightarrowG$ for the reflexive and transitive closure of 
$\arrG\tau$;
i.e., $K
\LongrightarrowG K'
 $ holds if $K $ can evolve into $K'$ by performing some
silent steps~-- possibly none.
We also write 
$K \ArrG \mu K' $ if   $K \LongrightarrowG \arrG\mu \LongrightarrowG K' $ (the
composition of the three 
relations); i.e.,  
$K \ArrG \mu K'$ holds if there are $K_1 $ and $K_2$ with 
$K \LongrightarrowG K_1$, $K_1 \arrG\mu K_2$ 
and $K_2 \LongrightarrowG K'$.
Then $ \ArcapG \mu  $  is  $\ArrG \mu $ if $\mu\neq \tau$, and  
$\LongrightarrowG $ if $\mu=\tau$. 

 \emph{Traces}, ranged over by $\trace$,  are finite (and possibly empty) sequences of
 visible actions.
If $\trace =\ell_1, \ldots, \ell_n$ ($n\geq 0$), then  $K \ArrG \trace K'$  holds if there
are $K_0, \ldots, K_n$ with $K_0 = K$, $K_n =K'$, and $K_i \ArrG{\ell_{i+1}} K_{i+1}$  for
$0\leq i < n$; and    $K \ArrG \trace $ if there is $K' $ with $K \ArrG \trace  K'$. 

In an action $\inp a\tilb$ or $\bout a\tilb$ we call $a$ the \emph{subject} and
$\tilb$ the object; similarly in a trace $s$ a name $a$ occurs in \emph{subject position}
(resp.\ in \emph{object position})
if there is an action in $s$ whose subject is $a$ (resp.\ whose object includes $a$). 

Two states $K_1, K_2 $ of the LTS   are \emph{trace equivalent}, written 
 $K_1 \TEG K_2$, if    ($K_1 \ArrG \trace $  iff 
$K_2 \ArrG \trace  $), for all $\trace$. 
 
Similarly,  \emph{bisimilarity}, written $\wbG $,
 is the largest symmetric relation on the state of the LTS such that 
whenever $K_1 \wbG K_2$ then 
\begin{itemize}
\item $K_1 \ArrG \mu  K_1'$ implies there is $K_2'$ with 
$K_2 \ArcapG \mu  K_2'$ and $K'_1 \wbG K'_2$.
\end{itemize} 
For instance, in the \piI LTS   $\arrPI \mu $ of Section~\ref{ss:piI}, 
$P \TEPI Q$ means that the \piI  processes $P$ and $Q$ are trace equivalent, and 
$P \wbPI Q$ means that they are bisimilar.

\begin{remark}[bound names]
\label{r:bn} 
%% NEW 
In an action $\inp a \tilb$ or $\bout a \tilb$ or $(\tilb)$, 
name $a$ is \emph{free} whereas $\tilb$ are \emph{bound}; 
the free and bound names of a trace are defined accordingly.


Throughout the paper, 
%equality of observations is taken up to renaming of the bound names. 
%In other terms,  
in any statement (concerning OGS or \piI), the bound names of an action or of a trace that
appear in the statement are supposed to be all \emph{fresh}; i.e.,  
all distinct from each other and from the free names of the objects in the statement. 
For instance, in a  trace $s= s_1, \mu, s_2$, a  bound name of $\mu$ should not appear in the
preceding prefix  $s_1$
(neither bound nor free), should not have other occurrences  in $\mu$, 
and may only appear free in the following suffix  $s_2$; moreover, if  
$s$ is the trace 
in the definition above  of trace equivalence,  any  bound names  of
$s$ should not appear free in the initial states $K_1,K_2$.  
Similarly, in the following definition of bisimilarity, the bound names of the action $\mu$ do
not appear free in the states $K_1,K_2$. 
In the same way, the convention applies to statements about operational correspondence between an
OGS state $K_1$ and a \piI state $K_2$ (the bound names of actions or traces mentioned in
the statement do not appear free in  $K_1,K_2$).
% An action $\mu$ or a trace $\trace$ is \emph{well-formed} if all its bound names  are
% distinct; and it is  \emph{well-formed for $K_1,K_2$} if it is well-formed and moreover
% all  bound names are distinct from the free names  in $K_1$, $K_2$.
\qed\end{remark}
